\section{The Polynomial Commitment Scheme}\label{sec:pcs}

We now use the fold dictionary to build a multilinear polynomial commitment scheme for tables, in the sense of \cref{def:pcs}.
The prover encodes a table $f$ as the Reed--Solomon codeword of $U_f = \sum_b f(b) K_b$.
An evaluation claim $\tilde f(z) = v$ is proved by a sumcheck (\cref{sec:proofs:sumcheck}) whose round challenges are also the fold challenges, exactly as in BaseFold~\cite{ZCF24}.
By \cref{cor:codeword-fold}, no change of representation is needed: the folds of the committed word track the partially restricted table that the sumcheck prover already holds.
We present the scheme as a public-coin IOP (\cref{def:iop}); its compilation into a non-interactive scheme is discussed in \cref{sec:proofs:bcs,sec:sound:pq}.

\subsection{Parameters and encoding}\label{sec:pcs:params}

\paragraph{Fields.}
The smooth domain lies in the base field $\Fq$, while tables, words, the evaluation point and the challenges live in a finite field $\F \supseteq \Fq$, the \emph{challenge field}.
Taking $\F = \Fq$ gives the plain scheme; taking $\F = \K$, an extension of $\Fq$, gives the scheme with extension-field challenges used in \cref{sec:impl}, in which honest tables have entries in $\Fq$.
All results of \cref{sec:pcs,sec:sound} hold for every such $\F$, with $|\F|$ in the error bounds.

\paragraph{Parameters.}
\begin{itemize}
  \item the number of variables $n \ge 1$, and $N = 2^n$;
  \item a rate $\rho = 2^{-R}$ with $R \ge 1$, and a smooth domain $L \le \Fq^\times$ of order $M = N/\rho = 2^{n+R}$ (\cref{def:domain}; it exists if and only if $2^{n+R} \mid q-1$, by \cref{lem:power});
  \item a number of folding rounds $\ell \in \iv{0}{n}$;
  \item a number of queries $\kappa \ge 1$.
\end{itemize}
For $j \in \iv{0}{\ell}$ let $L_j = L^{2^j}$, the smooth domain of order $M_j = M/2^j \ge 2^R \ge 2$ (\cref{lem:power}), let $N_j = N/2^j$, and let $\Cc_j = \RS_\F[L_j, N_j]$.
Every code $\Cc_j$ has the same rate $N_j/M_j = \rho$, and $L_{j+1} = (L_j)^2$.

\begin{definition}[Encoding]\label{def:encoding}
The \emph{encoding} of a table $f : \iv{0}{N-1} \to \F$ is the word
\[
  \Enc(f) \;=\; \ev_L(U_f) \;\in\; \Cc_0,
  \qquad
  U_f \;=\; \sum_{b=0}^{N-1} f(b)\, K_b \;\in\; \F[X]_{<N} .
\]
\end{definition}

\begin{lemma}[Encoding]\label{lem:encoding}
$\Enc$ is an $\F$-linear bijection from the tables over $\F$ onto $\Cc_0$, and for $c \in \Cc_0$, the table $f$ with $\Enc(f) = c$ is the table of kernel coordinates of $P_c$.
\end{lemma}

\begin{proof}
$f \mapsto U_f$ is a linear bijection onto $\F[X]_{<N}$ by \cref{thm:basis}, and $\ev_L$ is a linear bijection from $\F[X]_{<N}$ onto $\Cc_0$ by \cref{lem:rs-params}(1), with inverse $c \mapsto P_c$.
\end{proof}

$\Enc$ is computed by the transform of \cref{cor:transforms} ($\tfrac{n}{2}N$ additions), followed by one number-theoretic transform (NTT, the fast evaluation of a polynomial on a smooth domain) of size $M$ on the zero-padded coefficient vector.

\subsection{The evaluation protocol}\label{sec:pcs:protocol}

The commitment to $f$ is the word $w_0 = \Enc(f)$.
The evaluation protocol $\Pieval$ proves the claim $\tilde f(z) = v$ for a public point $z = (z_1,\dots,z_n) \in \F^n$ and a public value $v \in \F$.
As in \cref{def:pcs}, its instance is $(w_0; z, v)$, and the verifier has oracle access to $w_0$.

\paragraph{Round structure.}
$\Pieval$ is a public-coin IOP with $\ell+1$ rounds.
\begin{itemize}
  \item In round $j \in \iv{1}{\ell}$, the prover sends a \emph{round polynomial} $s_j \in \F[X]_{<3}$, together with, if $j \ge 2$, an oracle $w_{j-1} \in \F^{L_{j-1}}$; the verifier replies with a uniformly random challenge $r_j \in \F$.
  \item In round $\ell+1$, the prover sends a \emph{final table} $g : \iv{0}{N_\ell-1} \to \F$ in the clear; the verifier replies with $\kappa$ independent uniformly random \emph{query points} $\xi^{(1)}_0,\dots,\xi^{(\kappa)}_0 \in L$, and then decides.
\end{itemize}
Thus the challenge sets are $\mathsf{Ch}_j = \F$ for $j \in \iv{1}{\ell}$ and $\mathsf{Ch}_{\ell+1} = L^{\times\kappa}$, the set of $\kappa$-tuples of elements of $L$ (a Cartesian power, not the image of a power map).
A round polynomial is transmitted as its values at $0$, $1$ and $2$, which are distinct because the characteristic is odd. The map $\F[X]_{<3} \to \F^3$, $s \mapsto (s(0), s(1), s(2))$, is injective by \cref{lem:roots}, hence bijective since both spaces have dimension $3$; so every triple sent by a prover identifies a unique $s_j \in \F[X]_{<3}$.
As in \cref{sec:proofs:iop}, the checks below are predicates on the full transcript.

\paragraph{Honest prover.}
The honest prover keeps two tables over $\F$, initialised to
\[
  f_0 = f \qquad\text{and}\qquad \chi_0 = \bigl(b \mapsto \eq(b,z)\bigr), \qquad b \in \iv{0}{N-1},
\]
and after round $j$ replaces them by their restrictions (\cref{def:restriction}) at $r_j$: $f_j = \res_{r_j}(f_{j-1})$ and $\chi_j = \res_{r_j}(\chi_{j-1})$, tables on $\iv{0}{N_j-1}$.
In round $j$ it sends
\begin{equation}\label{eq:honest-sj}
  s_j(X) \;=\; \sum_{b=0}^{N_j - 1}
      \bigl((1-X)\, f_{j-1}(2b) + X f_{j-1}(2b+1)\bigr)\,
      \bigl((1-X)\, \chi_{j-1}(2b) + X \chi_{j-1}(2b+1)\bigr),
\end{equation}
and, if $j \ge 2$, the oracle $w_{j-1} = \fold_{r_{j-1}}(w_{j-2})$ (\cref{def:fold-word}).
In round $\ell+1$ it sends $g = f_\ell$.

\paragraph{Verifier.}
Set $v_0 = v$ and $v_j = s_j(r_j)$ for $j \in \iv{1}{\ell}$.
Let $G = \sum_{b=0}^{N_\ell - 1} g(b)\, K^{(n-\ell)}_b \in \F[X]_{<N_\ell}$ and $w_\ell = \ev_{L_\ell}(G)$; the verifier computes values of $w_\ell$ itself.
The verifier accepts if and only if all of the following checks pass.
\begin{enumerate}
  \item \emph{Sumcheck.} $s_j(0) + s_j(1) = v_{j-1}$ for every $j \in \iv{1}{\ell}$.
  \item \emph{Closure.}
  \[
    v_\ell \;=\; \Bigl(\prod_{j=1}^{\ell} \eq_1(r_j, z_j)\Bigr)\cdot \tilde g\bigl(z_{\ell+1},\dots,z_{n}\bigr).
  \]
  \item \emph{Queries.} For every $i \in \iv{1}{\kappa}$, with $\xi_0 = \xi^{(i)}_0$ and $\xi_j = \xi_0^{2^j}$ for $j \in \iv{1}{\ell}$:
  \begin{itemize}
    \item if $\ell \ge 1$: for every $j \in \iv{1}{\ell}$, the value $\fold_{r_j}(w_{j-1})(\xi_j)$, computed by \eqref{eq:fold-explicit} from the two queried values $w_{j-1}(\pm\xi_{j-1})$, equals $w_j(\xi_j)$ (a query to the oracle $w_j$ if $j < \ell$, and $G(\xi_\ell)$ if $j = \ell$);
    \item if $\ell = 0$: $w_0(\xi_0) = G(\xi_0)$.
  \end{itemize}
\end{enumerate}

\begin{remark}[Stopping parameter]\label{rem:stop}
With $\ell = n$, the final table is the single value $\tilde f(r_1,\dots,r_n)$ and $G$ is a constant; with $\ell = 0$, the prover sends $f$ in the clear.
Intermediate values trade the $N_\ell = 2^{n-\ell}$ field elements of the final message against $\ell$ folding rounds.
\end{remark}

\subsection{Completeness}\label{sec:pcs:complete}

Recall from \cref{sec:prelim:bits} the notation $(r_{\le j}, b) = (r_1,\dots,r_j,b_1,\dots,b_{n-j})$ for an index $b$ with $n-j$ bits.

\begin{lemma}[Honest prover]\label{lem:honest}
Let $f$ be a table over $\F$, $z \in \F^n$, and $r_1,\dots,r_\ell \in \F$, and let $f_j$, $\chi_j$ and $s_j$ be computed by the honest prover. Then, for $j \in \iv{0}{\ell}$ and $b \in \iv{0}{N_j-1}$:
\begin{enumerate}
  \item $f_j(b) = \tilde f(r_{\le j}, b)$ and $\chi_j(b) = \eq\bigl((r_{\le j}, b), z\bigr)$;
  \item for $j \ge 1$, $s_j$ is the honest round polynomial $s^*_j$ of \cref{sec:proofs:sumcheck} for $F(x) = \tilde f(x)\,\eq(x,z)$, and $s_j \in \F[X]_{<3}$;
  \item for $j \ge 1$, $s_j(r_j) = \sum_b f_j(b)\,\chi_j(b)$, and $s_j(0) + s_j(1) = \sum_b f_{j-1}(b)\,\chi_{j-1}(b)$;
  \item $\sum_b f_0(b)\chi_0(b) = \tilde f(z)$ and, for every $j$,
  \[
    \sum_{b=0}^{N_j-1} f_j(b)\,\chi_j(b) \;=\; \Bigl(\prod_{i=1}^{j}\eq_1(r_i,z_i)\Bigr)\cdot \widetilde{f_j}(z_{j+1},\dots,z_n).
  \]
\end{enumerate}
\end{lemma}

\begin{proof}
(1) \Cref{lem:restriction}(2), and \cref{cor:eq-mle} for $\chi_0$.
(2) Regard $f_{j-1}$ and $\chi_{j-1}$ as tables over the ring $\F[X] \supseteq \F$ (\cref{sec:prelim:ml}). The two factors in \eqref{eq:honest-sj} are then $\res_X(f_{j-1})(b)$ and $\res_X(\chi_{j-1})(b)$, and \cref{lem:restriction}(2) over $\F[X]$, with (1), gives
\[
  s_j(X) \;=\; \sum_{b=0}^{N_j-1} \tilde f(r_1,\dots,r_{j-1}, X, b_1,\dots,b_{n-j})\;\eq\bigl((r_1,\dots,r_{j-1},X,b_1,\dots,b_{n-j}),z\bigr) \;=\; s^*_j(X).
\]
Each summand of \eqref{eq:honest-sj} is a product of two polynomials of degree $\le 1$ in $X$, so $s_j \in \F[X]_{<3}$.
(3) By \cref{lem:natural} applied to the evaluation $X \mapsto r_j$, the factors become $f_j(b)$ and $\chi_j(b)$. Setting $X = 0$ and $X = 1$ in \eqref{eq:honest-sj} gives $\sum_b \bigl(f_{j-1}(2b)\chi_{j-1}(2b) + f_{j-1}(2b+1)\chi_{j-1}(2b+1)\bigr)$.
(4) The first identity is \cref{prop:mle}. For the second, (1) and \cref{lem:eq}(3) give
\[
  \chi_j(b) \;=\; \Bigl(\prod_{i=1}^{j}\eq_1(r_i,z_i)\Bigr)\,\eq\bigl(b,(z_{j+1},\dots,z_n)\bigr);
\]
conclude with \cref{prop:mle}.
\end{proof}

\begin{theorem}[Completeness]\label{thm:complete}
If $w_0 = \Enc(f)$, $v = \tilde f(z)$ and the prover is honest, the verifier accepts with probability $1$.
\end{theorem}

\begin{proof}
\emph{Sumcheck.} By \cref{lem:honest}(3,4), $v_0 = \tilde f(z) = \sum_b f_0(b)\chi_0(b)$ and $v_j = s_j(r_j) = \sum_b f_j(b)\chi_j(b)$, so $s_j(0) + s_j(1) = \sum_b f_{j-1}(b)\chi_{j-1}(b) = v_{j-1}$ for every $j \in \iv{1}{\ell}$.

\emph{Closure.} By \cref{lem:honest}(4) with $j = \ell$ and $g = f_\ell$, $v_\ell = \prod_{j}\eq_1(r_j,z_j) \cdot \tilde g(z_{\ell+1},\dots,z_n)$.

\emph{Queries.} If $\ell = 0$, then $g = f$, $G = U_f$ and $w_0 = \Enc(f) = \ev_L(G)$.
If $\ell \ge 1$, extend $(r_1,\dots,r_\ell)$ arbitrarily to $r \in \F^n$, and let $U_j$ be as in \cref{cor:iterated} for $U = U_f$; its kernel coordinates are $f_j$.
Since $L$ has order $2^{n+R}$ with $n+R \ge n$, \cref{cor:codeword-fold} shows that the successive folds of $w_0$ satisfy $\fold_{r_j}(\cdots\fold_{r_1}(w_0)) = \ev_{L_j}(U_j)$ for $j \in \iv{1}{\ell}$.
The honest oracles are these folds for $j < \ell$, and $U_\ell = \sum_b f_\ell(b) K^{(n-\ell)}_b = G$ because $g = f_\ell$.
Hence $\fold_{r_j}(w_{j-1}) = w_j$ on all of $L_j$ for every $j \in \iv{1}{\ell}$, and every check of every query passes.
\end{proof}

\subsection{Costs}\label{sec:pcs:costs}

Let $t_{\mathsf{H}}$ denote the cost of one hash evaluation. We count field operations up to constant factors.

\begin{itemize}
  \item \emph{Commit.} $\tfrac{n}{2}N$ additions (\cref{cor:transforms}), one NTT of size $M$, that is, $O(M\log M)$ operations, and a Merkle tree over the $M/2$ fibres of $w_0$ (one leaf per fibre), $O(M)\cdot t_{\mathsf{H}}$.
  \item \emph{Prover, evaluation.} Sumcheck: $O(N)$ in total, since each round is linear in the current table size and the sizes halve.
  Folds: $\fold_{r_j}(w_{j-1})$ costs $O(M_j)$ using precomputed inverses of $2\xi$ for $\xi \in L$; $O(M)$ in total.
  Merkle trees for $w_1,\dots,w_{\ell-1}$: $O(M)\cdot t_{\mathsf{H}}$ in total.
  \item \emph{Verifier.} Sumcheck: $O(\ell)$. Closure: $O(N_\ell)$. Queries: $O(\kappa \ell)$ field operations plus $O(\kappa \ell \log M)\cdot t_{\mathsf{H}}$ for authentication paths, plus $\kappa$ evaluations of $G$, $O(\kappa N_\ell)$.
  \item \emph{Proof size.} For $\ell \ge 1$: $3\ell$ field elements for the sumcheck, $N_\ell$ for the final table, $\ell-1$ Merkle roots besides the commitment, and $\kappa\ell$ opened leaves ($2\kappa \ell$ values) with their authentication paths. For $\ell = 0$: the table $g = f$ and $\kappa$ opened values of $w_0$.
\end{itemize}

\begin{remark}[Relation to BaseFold and DeepFold]\label{rem:basefold}
The interleaving of sumcheck rounds with fold rounds under shared challenges is the BaseFold paradigm~\cite{ZCF24}, and the asymptotic costs above match Reed--Solomon instantiations of BaseFold and DeepFold~\cite{GLHQTZ25}.
In particular, converting a table to coefficient form by \cref{lem:mobius} also costs $\tfrac{n}{2}N$ additions, so the kernel encoding offers no asymptotic advantage.
What differs is structural: by \cref{thm:fold-dictionary}, the committed word, the sumcheck table and the folds are three views of \emph{one} object, the restricted table $f_j$, and the consistency between them is an identity rather than a bookkeeping convention (compare \cref{rem:pitfall}).
\end{remark}

\begin{remark}[Batching and zero knowledge]\label{rem:batch}
Because $\Enc$ and $\fold_r$ are linear, several tables can be opened at the same point by running $\Pieval$ on a random linear combination of them, the verifier answering queries to the combined word from queries to each committed word.
We do not analyse batching in this paper, and zero knowledge is out of scope.
\end{remark}
